\chapter{Concluding Remarks}
\label{sec:conclusions}
% ============================================


% Springer Nature: each chapter should have a 10--15 line online abstract.
% This chapter abstract has not yet been written in the uploaded manuscript.
% [October 2026 round] Abstract written (blue).
\abstract{This closing chapter draws together the unified view of sensor fusion developed in the book. It restates the central thesis, namely that classical estimation and learned perception are layers of one architecture rather than competing paradigms, and that additional sensors improve a system only when their measurement models, uncertainties, and mutual dependencies are represented correctly. A synthesis table maps the recurring fusion tasks of the book (ego-localisation, object tracking, perception, occupancy, reliability assessment, anomaly detection, and cooperative fusion) to the classical tool, the learned component, and the hybrid interface that the book recommends for each. The chapter then records the key take-aways of each part of the book, from the foundations and classical techniques through learning-based and hybrid fusion, the autonomous-driving application, multi-agent fusion, and the case studies beyond driving, and closes with guidance on how to use the book, its end-of-chapter problems, and the companion repository of reproducible Python examples in teaching, research, and engineering practice.}
\suppressfloats[t]

This book has presented a tutorial and survey of sensor fusion for autonomous driving vehicles. By focusing on autonomous driving as the central application domain, the book has connected classical fusion methods, learning-based multimodal perception, reliability-aware reasoning, practical implementation, and safety-critical evaluation. Classical methods such as Bayesian filtering, statistical fusion, evidence theory, and optimisation remain essential for localisation, tracking, uncertainty management, and sensor-health reasoning. Learning-based methods, particularly deep multimodal and bird's-eye-view fusion, have become central to perception tasks such as object detection, segmentation, tracking, and occupancy estimation.

Set-membership and zonotopic estimation add a conservative bounded-error layer for consistency checking, complementing probabilistic and learning-based fusion when explicit guarantees are needed.

A central theme of this book is that robust autonomous driving requires more than combining sensor outputs for improved benchmark accuracy. Fusion systems must also handle calibration errors, asynchronous data, adverse weather, sensor degradation, out-of-distribution scenes, computational constraints, and uncertainty propagation. The interaction between fusion outputs and downstream prediction, planning, and control means that safety, latency, interpretability, and fault tolerance must be considered throughout the design process.

% % Checked and approved.
The main lesson of the book is that additional sensors improve a system only when their measurement models, uncertainties, and mutual dependencies are represented correctly: an ignored correlation, a miscalibrated covariance, or an unmodelled failure mode can make a fused estimate overconfident, and an overconfident estimate is more dangerous than an uncertain one because the modules downstream trust it. A consequence for the architecture is that learned perception and classical estimation can be used as layers of the same system, the classical layer propagating and checking the uncertainty of the learned outputs before they are acted upon.

\subsection*{The Unified View in One Table}

The architectural corollary can be made concrete. Every fusion task treated in the book admits the same decomposition into a classical tool that represents and propagates uncertainty, a learned component that supplies recognition or prior knowledge that no hand-designed model can match, and an interface through which the two exchange information in a form the other can audit. Table~\ref{tab:unified-view} collects the recommended decomposition for the seven recurring tasks. Reading the table column by column shows why the view is unified: the interfaces of the third column are drawn from a small set of mechanisms (a learned covariance used as the measurement-noise matrix, a learned confidence used as a discount factor, innovation and conflict statistics used as monitors, and covariance intersection used whenever correlations are unknown), and each of them is defined in Chapter~\ref{sec:classicalSF}, instantiated in Chapters~\ref{sec:learning} and~\ref{sec:application}, extended to several agents in Chapter~\ref{sec:MAS}, and reused in the domains of Chapter~\ref{sec:applications}.

% [October 2026 round] Synthesis table requested by the editor; the \color line is only needed while \revisionmarkuptrue.
\begin{table}[!t]
\caption{The unified view: for each recurring fusion task, the classical tool, the learned component, and the hybrid interface recommended in this book.}
\label{tab:unified-view}
\footnotesize
\renewcommand{\arraystretch}{1.25}% row spacing (layout only)
\setlength{\tabcolsep}{3pt}
\begin{tabular}{>{\raggedright\arraybackslash}p{0.16\textwidth} >{\raggedright\arraybackslash}p{0.26\textwidth} >{\raggedright\arraybackslash}p{0.24\textwidth} >{\raggedright\arraybackslash}p{0.27\textwidth}}
\toprule
Fusion task & Classical tool & Learned component & Hybrid interface \\
\midrule
Ego-localisation & EKF/UKF and factor-graph smoothing of GNSS, IMU, and odometry (Chapters~\ref{sec:classicalSF}, \ref{sec:learning}); zonotopic enclosures for integrity & Learned visual features, place recognition, and loop-closure detection & Learned correspondences enter as measurements with covariance; innovation tests and enclosure containment audit the residuals \\[2pt]
Object tracking & Kalman tracking, gating \eqref{eq:mahalanobis_gate}, JPDA, track-to-track fusion with covariance intersection & Neural multimodal 3D detector with learned covariance & Learned covariance used as $\mathbf{R}$; filter residuals fed back as a calibration signal (Fig.~\ref{fig:hybrid-pipeline}) \\[2pt]
Perception (detection, segmentation) & Evidence-theoretic combination of per-modality outputs; geometric priors & Deep early, feature, and late fusion; BEV and transformer architectures & Per-modality outputs mapped to belief masses with discounting; conflict mass $K$ as a monitor \\[2pt]
Occupancy & Probabilistic grids and log-odds updates \eqref{eq:octomap_update} with inverse sensor models & Semantic occupancy networks & Network outputs treated as per-cell measurements whose weight depends on sensor reliability \\[2pt]
Reliability assessment & MCDM/AHP weights \eqref{eq:mcdm-score}, evidence discounting, NIS monitoring \eqref{eq:nis} & Learned confidence, uncertainty, and out-of-distribution scores & Context selects MCDM weights; learned confidence becomes a discount factor; innovation tests audit both \\[2pt]
Anomaly detection & Innovation tests, set-membership consistency, fuzzy rules & Learned anomaly scores on AIS tracks or distributed acoustic data (Chapter~\ref{sec:applications}) & Learned scores enter evidence combination as one source; set-based checks detect incompatibility when the assumed error bounds and models hold \\[2pt]
Cooperative fusion & Consensus \eqref{eq:consensus-filter}, covariance intersection \eqref{eq:ci-est}, decentralised filters (Chapter~\ref{sec:MAS}) & Feature-level sharing networks; federated learning & Exchanged tracks or features carry covariance and provenance; intersection for unknown correlations; discounting by trust \\
\bottomrule
\end{tabular}
\end{table}

\subsection*{Key Take-Aways by Part of the Book}

\textbf{Foundations and classical techniques (Chapters~\ref{sec:Principles} and~\ref{sec:classicalSF}).} Fusion is an inverse problem whose solution is only as good as its measurement models. Least squares, maximum a posteriori estimation, and Bayesian filtering are one family, differing in what they assume about priors and dynamics; the innovation sequence is the single most useful diagnostic that this family provides, because it tests the model against the data without ground truth. Three complementary representations of uncertainty should be kept in the toolbox, since each answers a question that the others cannot: covariances for efficient recursive estimation, bounded sets and zonotopes for guarantees when noise distributions are unknown but bounded, and belief functions for ignorance and conflict. Track-level fusion with covariance intersection and MCDM-based weighting round out the classical layer with the two decisions every practical system must take, how to combine estimates whose correlation is unknown and how much to trust each source in a given context.

\textbf{Learning-based and hybrid fusion (Chapter~\ref{sec:learning}).} Deep multimodal networks have redefined what is achievable in perception, and the design space of early, feature, late, bird's-eye-view, and occupancy-based fusion is now well charted. For the view of this book, these networks are most useful when they output an uncertainty together with each label, and when that uncertainty is consumed by a classical layer that can check it; in the hybrid pipeline of the chapter each paradigm is used for the part of the problem it handles best. The strongest challenge to this view comes from end-to-end architectures that learn perception, prediction, and planning jointly (Sect.~\ref{subsec:oc-foundation}); whether a classical safety layer can be retained when the fusion function becomes implicit is, in our view, the central open question for the next generation of fusion systems.

\textbf{The autonomous-driving application (Chapter~\ref{sec:application}).} A complete stack makes the abstract tools concrete: the sensor suite, the fusion levels, the core tasks, the datasets and metrics, and the implementation and validation infrastructure. Two lessons recur. Evaluation must go beyond average-case accuracy to robustness, calibration, latency, and the consequences of fusion errors for planning; and timing, calibration, and coordinate frames must be treated as part of the measurement model, since a timestamp offset or a stale extrinsic calibration is indistinguishable from a sensor fault.

\textbf{Multi-agent fusion (Chapter~\ref{sec:MAS}).} Sharing information across vehicles and infrastructure extends every classical mechanism rather than replacing it. Consensus and decentralised filtering distribute the estimation; covariance intersection handles the unknown correlations that any network with multiple paths creates; discounting handles senders of unknown reliability. The failure modes specific to this setting, double counting, rumour propagation, inconsistent time bases, and spoofing, all stem from an incorrect description of the uncertainty of the exchanged information, which is why the same tools address them.

\textbf{Beyond driving (Chapter~\ref{sec:applications}).} The maritime and submarine-cable case studies show that the unified view applies beyond road vehicles: the same evidence-theoretic combination, reliability weighting, consistency checking, and hybrid use of learned anomaly scores apply to AIS tracks and distributed acoustic sensing, with different sensors, time scales, and consequences of error. Readers working in other safety-critical domains should find in these chapters a template rather than a special case, and Chapter~\ref{sec:open-challenges} shows that the open problems, too, are shared across domains.

\subsection*{How to Use This Book}

The chapters can be read in sequence as a course, or selectively: Chapters~\ref{sec:Principles} and~\ref{sec:classicalSF} form a self-contained introduction to estimation-based fusion, Chapters~\ref{sec:learning} and~\ref{sec:application} a self-contained introduction to learning-based perception and its deployment, and Chapters~\ref{sec:MAS} to~\ref{sec:open-challenges} are intended for readers who already know one of the two traditions and want to see how the other completes it. Each technical chapter ends with problems that mix derivations, design questions, and small computations; model solutions are collected at the end of the book. The computational problems are supported by the companion repository, which contains reproducible Python examples organised chapter by chapter (\repopath{companion\_\allowbreak repo/\allowbreak examples} and \repopath{companion\_\allowbreak repo/\allowbreak exercises/\allowbreak chNN}), the scripts that generate the data behind the figures, and a notebook version of every script (\repopath{*.ipynb}) so that the examples can be run and modified interactively; the data files needed to compile the book are shipped with it, so that the manuscript does not depend on the repository, while the numerical figures and computational tables can be regenerated using the stated data and settings. We recommend running the examples alongside the text, modifying the noise models, reliability weights, and timestamp offsets, and observing how the consistency checks of the classical layer react; this is the most direct way to see how the description of uncertainty drives the behaviour of a fusion system.

Future research is likely to move towards hybrid architectures that combine the representational power of deep learning with the structure, uncertainty handling, and interpretability of classical estimation and decision-making methods. Progress will also depend on improved datasets, more realistic robustness benchmarks, efficient embedded deployment, and rigorous verification and validation. Sensor fusion will therefore remain a core research and engineering challenge in the development of safe, reliable, and scalable autonomous driving vehicles. %
In our view, the most useful near-term contributions will come from the careful integration of the two traditions reviewed here, and from shared benchmarks that reward calibrated uncertainty and graceful degradation in addition to average-case accuracy.
